\begin{ouverture}[espace]{18}{Solides, assemblages et volumes}{Partie 3 · Espace et géométrie}
{Solides usuels · Vues · Cubes cachés · Centimètre cube}
\objectif{Reconnaître les solides usuels et utiliser leur vocabulaire.}
\objectif{Passer d’un assemblage de cubes à ses vues, et inversement.}
\objectif{Déterminer et comparer des volumes en centimètres cubes.}
\prerequis{Solides et patrons du cours moyen ; dénombrement organisé.}
\end{ouverture}
\pagestyle{otMain}\essentiel\fontsize{11}{15}\selectfont
\phantomsection\label{cours-18-1}
\essSection{1}{Reconnaître et décrire les solides usuels}
\begin{definition}Le cube et le pavé droit ont 6 faces rectangulaires. Un prisme droit a deux bases polygonales identiques et des faces latérales rectangulaires. Une pyramide a une base polygonale et des faces triangulaires qui se rejoignent en un sommet. Le cylindre, le cône et la boule ont des surfaces courbes.\end{definition}
\label{fig-cours-18-1-cours-solides}\figureGeo{cours-solides}

\begin{propriete}Les faces d’un polyèdre sont planes ; deux faces se rencontrent sur une arête ; les arêtes se rejoignent en des sommets. Le pavé droit a 6 faces, 12 arêtes et 8 sommets.\end{propriete}
\essSection{2}{Comprendre les vues et les perspectives}
\begin{definition}Une perspective représente un solide sur une feuille ; les arêtes cachées y sont en pointillés. Une vue de face, de dessus ou de côté montre ce que l’on voit en regardant le solide dans une seule direction.\end{definition}
\clearpage
\phantomsection\label{cours-18-2}
\essSection{3}{Dénombrer les cubes visibles et cachés}
\label{fig-cours-18-2-cours-assemblage-vues}\figureGeo{cours-assemblage-vues}

\begin{savoirfaire}[Compter par étages]On compte étage par étage, sans oublier les cubes cachés qui portent ceux du dessus. Sur la vue de dessus, on peut écrire dans chaque case le nombre de cubes empilés : le total des nombres donne le nombre de cubes.\end{savoirfaire}
\essSection{4}{Mesurer un volume en centimètres cubes}
\begin{definition}Le volume d’un solide mesure la place qu’il occupe dans l’espace. 1 cm$^3$ est le volume d’un cube de 1 cm d’arête. Un assemblage de 14 cubes de 1 cm d’arête a un volume de 14 cm$^3$.\end{definition}
\begin{attention}Deux solides de formes différentes peuvent avoir le même volume : il suffit qu’ils contiennent le même nombre de cubes unités.\end{attention}

\clearpage\rubrique{Découverte}
\begin{blocLivre}{ex-18-1}\begin{exercice}[1]{Nommer les solides}
Nommer chacun des six solides numérotés.
\label{fig-ex-18-1}\figureGeo{sol-1}
\end{exercice}\end{blocLivre}
\begin{blocLivre}{ex-18-2}\begin{exercice}[1]{Faces, arêtes, sommets}
Compter les faces, les arêtes et les sommets du pavé droit, puis de la pyramide à base carrée.
\label{fig-ex-18-2}\figureGeo{sol-2}
\end{exercice}\end{blocLivre}
\begin{blocLivre}{ex-18-3}\begin{exercice}[1]{Un petit cube de cubes}
Le grand cube est formé de cubes de 1 cm d’arête. Combien de petits cubes le composent ? Quel est son volume ?
\label{fig-ex-18-3}\figureGeo{sol-3}
\end{exercice}\end{blocLivre}
\begin{blocLivre}{ex-18-4}\begin{exercice}[1]{La bonne vue de dessus}
Parmi les vues 1, 2 et 3, laquelle est la vue de dessus de l’assemblage ?
\label{fig-ex-18-4}\figureGeo{sol-4}
\end{exercice}\end{blocLivre}
\begin{blocLivre}{ex-18-5}\begin{exercice}[1]{L’escalier}
L’escalier est formé de cubes de 1 cm d’arête. Quel est son volume ?
\label{fig-ex-18-5}\figureGeo{sol-5}
\end{exercice}\end{blocLivre}

\rubrique{Entraînement}
\begin{blocLivre}{ex-18-6}\begin{exercice}[2]{Les cubes cachés}
L’assemblage est formé de trois étages pleins : un carré de 3 sur 3, un carré de 2 sur 2, puis un cube. Combien de cubes compte-t-il ? Combien sont cachés ?
\label{fig-ex-18-6}\figureGeo{sol-6}
\end{exercice}\end{blocLivre}
\begin{blocLivre}{ex-18-7}\begin{exercice}[2]{Dessiner les vues}
Dessiner sur quadrillage la vue de face, la vue de dessus et la vue de droite de l’assemblage.
\label{fig-ex-18-7}\figureGeo{sol-7}
\end{exercice}\end{blocLivre}
\begin{blocLivre}{ex-18-8}\begin{exercice}[2]{Le plan coté}
La vue de dessus indique dans chaque case le nombre de cubes empilés. Combien de cubes compte l’assemblage ? Quel est son volume ?
\label{fig-ex-18-8}\figureGeo{sol-8}
\end{exercice}\end{blocLivre}
\begin{blocLivre}{ex-18-9}\begin{exercice}[2]{Comparer deux volumes}
Les deux assemblages sont formés de cubes identiques. Lequel a le plus grand volume ?
\label{fig-ex-18-9}\figureGeo{sol-9}
\end{exercice}\end{blocLivre}
\begin{blocLivre}{ex-18-10}\begin{exercice}[2]{Compléter en grand cube}
L’assemblage compte 20 cubes. Combien faut-il en ajouter pour obtenir un cube de 3 cubes d’arête ?
\label{fig-ex-18-10}\figureGeo{sol-10}
\end{exercice}\end{blocLivre}

\rubrique{Problèmes}
\begin{blocLivre}{ex-18-11}\begin{exercice}[3]{La boîte de sucres}
Des morceaux de sucre de 1 cm$^3$ remplissent une boîte : 4 sucres en longueur, 3 en largeur et 2 en hauteur. Combien de sucres contient la boîte ? Quel est son volume ?
\label{fig-ex-18-11}\figureGeo{sol-11}
\end{exercice}\end{blocLivre}
\begin{blocLivre}{ex-18-12}\begin{exercice}[3]{Le mur percé}
Un mur est formé de cubes de 1 dm d’arête : 6 cubes de long, 4 de haut et 1 d’épaisseur, avec une fenêtre de 2 cubes. Combien de cubes compte le mur ?
\label{fig-ex-18-12}\figureGeo{sol-12}
\end{exercice}\end{blocLivre}
\begin{blocLivre}{ex-18-13}\begin{exercice}[3]{Le grand cube peint}
Un cube de 3 cubes d’arête est peint sur toutes ses faces, puis séparé en 27 petits cubes. Combien de petits cubes ont 3 faces peintes ? 2 faces ? 1 face ? Aucune ?
\label{fig-ex-18-13}\figureGeo{sol-13}
\end{exercice}\end{blocLivre}
\begin{blocLivre}{ex-18-14}\begin{exercice}[3]{Des vues à l’assemblage}
On donne la vue de dessus, la vue de face et la vue de droite d’un assemblage de cubes posé sur une table. Combien de cubes le composent ?
\label{fig-ex-18-14}\figureGeo{sol-14}
\end{exercice}\end{blocLivre}
\begin{blocLivre}{ex-18-15}\begin{exercice}[3]{Même volume ?}
Chloé affirme : « Les deux assemblages n’ont pas la même forme, donc ils n’ont pas le même volume. » A-t-elle raison ?
\label{fig-ex-18-15}\figureGeo{sol-15}
\end{exercice}\end{blocLivre}

\clearpage
